\cleardoublepage
\chapter*{后记}
\addcontentsline{toc}{chapter}{后记}
\thispagestyle{empty}

一本数学书通常会在数学本身尚未结束的地方结束。

最后一个定理或许已经证明，最后一个例子或许已经讨论，所有记号也暂时抵达了一个停顿点。然而，这些都不构成真正的结论。本书中引入的每一个定义都可以继续抽象，每一个定理都属于一个更大的理论，而每一个证明都会引出新的问题：哪些假设不可缺少，哪些论证可以推广，还有哪些结构尚未被发现。

因此，这些笔记的终点不应被理解为一段数学旅程的完成。它只标记了某一条具体路线被记录到这里。

这份手稿最初源于我整理所学知识的尝试。和许多初次接触大学数学的学生一样，我最早遇到的是一门门彼此分开的课程：分析研究极限与连续，线性代数研究向量与矩阵，抽象代数研究群、环与域，拓扑研究开集与连续映射，概率论研究随机变量与分布。每一门学科似乎都有自己的术语、技巧和论证标准。

它们共同的结构，是后来才逐渐显现的。

同一批问题不断以不同形式返回：我们正在研究什么对象？对象之间哪些关系真正重要？哪些变换保持相关结构？一个结论到底需要哪些假设？何时可以把局部性质推广到整体？何时代数描述对应于几何结构？何时收敛能够保持我们希望执行的运算？

一旦认识到这些问题，数学就不再像一组相互独立的领土，而更像一个由重复出现的思想构成、以不同语言表达的网络。

完备性最初作为数系与度量空间的性质出现，后来却控制极限、定点与解的存在；紧致性起初是一个拓扑条件，却一次次把局部信息转化为整体控制；对称性在代数中由群编码，在几何中由变换编码，在分析中则表现为不变性；线性最初出现在方程组中，随后成为把复杂对象分解为可处理部分的原则；测度扩展了初等的长度与面积概念，而概率又赋予同一结构一种新的解释。

因此，学习数学并不只是收集结论，更是学会在这些思维模式再次出现时辨认它们。

写作使这一课变得更加严格。阅读一个定理时，人很容易相信自己已经理解它；但解释其假设为什么自然、在不隐藏步骤的情况下重构证明，并把它放回周围理论中，则困难得多。数学表述会暴露“熟悉”与“理解”之间的差别。一个看似无害的句子可能隐藏着未定义的术语，一个“显然”的推论可能需要单独的引理，而一个标准证明也可能依赖许多页之前确立的约定。

从这个意义上说，写作不是学习完成后的记录，而是学习本身的一部分。

这份手稿的不完善也属于同一过程。各个学科的规模并不相等，单独一卷也不可能以相同深度处理它们。有些章节走近某个理论的前沿，另一些只建立继续学习所需要的语言；有些论证因为能够照亮一般方法而被纳入，另一些同样重要的结果则因需要另一整本书才能妥善处理而被省略。这里呈现的不是数学的最终地图，而是一名学生试图看清其地貌时所走过的路径。

这条路径必然还会改变。

曾经显得人为的定义，后来可能变得不可避免；曾经困难的证明，后来可能只是某种熟悉技巧的例子。反过来，看似已经解决的命题，也可能在另一门学科的视角下显露出更深的问题。数学成熟不会消除困难，它只会改变困难出现的层级。

初学者会问怎样证明一个定理。更有经验的学习者还会问：这个定理为什么这样表述？它的假设如何被发现？结论是否已经最优？究竟是什么更大的结构使这个证明成为可能？进步不仅由能够解决多少问题衡量，也由学会提出何种质量的问题衡量。

接受不完备本身，也需要一种纪律。

人很容易想把写作推迟到掌握所有细节之后，或者等每一章都达到同样精细的程度再分享笔记。但数学很少允许所谓“完全掌握”的状态。每一个答案都会扩大未知的边界。等待最终完成，往往意味着永远不会开始。

更诚实的目标不是完美，而是持续修正：把假设写得更清楚，用精确论证替代模糊直觉，为草率命题寻找反例，并在更好的结构显现时重写解释。严谨性不是一次取得后便会自动保持的属性，而是一种反复返回已写文字、追问它是否真正表达原意的习惯。

读者也应当以同样的态度对待这些页面。

不要把证明当成只能远观的纪念碑。把它们拆开，重新构造，寻找更短的论证与更强的结论，追问删去一个假设后会发生什么；画出符合定义的例子，也构造几乎符合定义的反例；比较同一思想的不同表述。只有当一个证明能够被重新创造、调整并质疑时，它才真正成为个人数学理解的一部分。

这些笔记也不应被视为终点。它们自然的延续存在于真正学习数学所依赖的书籍、课程、习题与讨论中。实分析通向泛函分析、偏微分方程与动力系统；代数通向表示论、交换代数与代数几何；拓扑通向流形、同伦论与几何；测度与概率通向随机过程、统计学与数学物理。而这些方向之间的边界，本身也越来越显得人为。

因此，序言中所说的数学大厦并不是一座已经竣工的摩天楼。它更像一座永久处于建设中的城市：一些地基极其古老，一些高塔刚刚建成；不同街区按不同需求发展，但连接它们的桥梁不断出现。当新的视角提出要求时，旧结构也会被加固、延伸或重建。

进入数学，并不只是从高处欣赏这座城市，而是学会在其中工作。

一切从小事开始：理解一个定义，检验一个例子，补全一个证明。久而久之，这些行动会发展为一种思维方式。精确不再只是形式义务，而成为智识诚实的一种形式；抽象不再是逃离具体现实，而成为发现许多具体现象共有结构的方法；证明也不再只是结论为真的证书，而是解释在给定假设下，一个命题为什么不能不如此成立。

这也许是我在写作这些笔记的过程中获得的最深体会。数学并不只由结论的确定性定义，它同样由思想从一步走向下一步时所保持的谨慎定义。

手稿在这里结束，但它的目的延伸到最后一页之外。如果它能够促使读者继续独立前进：质疑既有论证，寻找遥远学科之间的联系，并最终写出比这里更清楚的解释，那么它就已经完成了自己的任务。

序言面向一条尚未走过的道路；后记回望过去，才发现道路也改变了行路的人。

前方的一切，比已经记录下来的更加辽阔。

\begin{flushright}
李谨硕 \\
上海交通大学
\end{flushright}

\cleardoublepage
